\chapter{System Pipeline and Dataset Creation}

This chapter provides a detailed explanation of the complete system pipeline and the custom dataset creation process due to the lack of publicly available ISL datasets.

\section{Dataset Challenge and Custom Solution}

\subsection{Problem: Dataset Unavailability}

The primary challenge faced during project development was the \textbf{unavailability of publicly available Indian Sign Language (ISL) datasets}. Publicly available sign language datasets are typically:

\begin{itemize}
    \item Limited to American Sign Language (ASL) or other international sign languages
    \item Restricted by licensing agreements
    \item Not freely accessible for commercial or research use
    \item Insufficient for training robust models
    \item Not tailored to specific recognition requirements
\end{itemize}

ISL-specific challenges:
\begin{itemize}
    \item Linguistic differences from ASL and other sign languages
    \item Region-specific gesture variations
    \item Limited online resources for ISL
    \item Lack of standardized ISL gesture libraries
\end{itemize}

\subsection{Solution: Custom Dataset Creation}

To overcome this challenge, we implemented a comprehensive custom dataset creation pipeline within the Streamlit dashboard:

\subsubsection{Dataset Collection Process}

\begin{enumerate}
    \item \textbf{Gesture Definition}: Define specific gestures to recognize (hello, thumbsup, etc.)
    \item \textbf{Training Mode}: Use dashboard training interface to capture images
    \item \textbf{Data Collection}: Capture 100+ images per gesture from multiple angles
    \item \textbf{Data Augmentation}: Vary lighting, distance, and hand position
    \item \textbf{Labeling}: Automatic labeling through directory organization
    \item \textbf{Validation}: Manual verification of collected data
\end{enumerate}

\subsubsection{Custom Dataset Statistics}

\begin{table}[H]
\centering
\begin{tabular}{|c|c|c|c|}
\hline
\textbf{Gesture} & \textbf{Images} & \textbf{Collection Date} & \textbf{Resolution} \\
\hline
Hello & 100 & 2026-02-04 & 640x480 \\
\hline
Thumbs Up & 100 & 2026-02-04 & 640x480 \\
\hline
Total & 200 & - & - \\
\hline
\end{tabular}
\caption{Custom ISL Training Dataset Summary}
\end{table}

\subsubsection{Data Collection Methodology}

\paragraph{Environment Setup}
\begin{itemize}
    \item Varied lighting conditions (natural, artificial, mixed)
    \item Different backgrounds (solid color, complex patterns)
    \item Multiple distances from camera (30cm to 100cm)
    \item Various hand positions within frame
\end{itemize}

\paragraph{Image Characteristics}
\begin{itemize}
    \item Resolution: 640×480 pixels
    \item Format: JPEG
    \item Color space: BGR (as captured by OpenCV)
    \item Frame rate: 30 FPS capture
    \item Total variety: 200+ unique samples
\end{itemize}

\section{Complete System Pipeline}

The system pipeline consists of 7 major stages:

\subsection{Stage 1: Video Input and Preprocessing}

\paragraph{Input Layer}
\begin{itemize}
    \item Capture from webcam (can also process video files)
    \item Resolution: 640×480 pixels at 30 FPS
    \item Color format: BGR (OpenCV standard)
\end{itemize}

\paragraph{Preprocessing}
\begin{lstlisting}[language=Python]
# Flip image for mirror effect
image = cv.flip(image, 1)

# Convert BGR to RGB for MediaPipe
image_rgb = cv.cvtColor(image, cv.COLOR_BGR2RGB)

# Create MediaPipe image
mp_image = mp.Image(image_format=mp.ImageFormat.SRGB, data=image_rgb)
\end{lstlisting}

\paragraph{Output}
\begin{itemize}
    \item Normalized RGB image (0-1 range)
    \item Timestamp in milliseconds for video mode
\end{itemize}

\subsection{Stage 2: Hand Landmark Detection (MediaPipe)}

\paragraph{Hand Detection Configuration}
\begin{lstlisting}[language=Python]
options = vision.HandLandmarkerOptions(
    base_options=base_options,
    num_hands=2,                              # Support dual-hand
    min_hand_detection_confidence=0.7,        # Detection threshold
    min_tracking_confidence=0.5,              # Tracking threshold
    running_mode=vision.RunningMode.VIDEO
)
hands = vision.HandLandmarker.create_from_options(options)
\end{lstlisting}

\paragraph{Landmark Structure}

MediaPipe provides 21 hand landmarks per hand:

\begin{table}[H]
\centering
\begin{tabular}{|c|c|c|}
\hline
\textbf{Index} & \textbf{Landmark} & \textbf{Position} \\
\hline
0 & Wrist & Base of hand \\
\hline
1-4 & Thumb & Base, middle, interphalangeal, tip \\
\hline
5-8 & Index & Base, PIP, DIP, tip \\
\hline
9-12 & Middle & Base, PIP, DIP, tip \\
\hline
13-16 & Ring & Base, PIP, DIP, tip \\
\hline
17-20 & Pinky & Base, PIP, DIP, tip \\
\hline
\end{tabular}
\caption{MediaPipe Hand Landmarks (21 points per hand)}
\end{table}

\paragraph{Output}
\begin{itemize}
    \item 21 landmarks per hand (x, y, z coordinates)
    \item Confidence scores (0-1)
    \item Handedness (LEFT/RIGHT)
    \item Detection results for multiple hands
\end{itemize}

\subsection{Stage 3: Landmark Preprocessing}

\paragraph{Normalization Process}

The raw landmarks require preprocessing before classification:

\begin{lstlisting}[language=Python]
def calc_landmark_list(image, landmarks):
    """Convert normalized coordinates to pixel values"""
    image_width, image_height = image.shape[1], image.shape[0]
    landmark_point = []
    for landmark in landmarks.landmark:
        landmark_x = min(int(landmark.x * image_width), 
                        image_width - 1)
        landmark_y = min(int(landmark.y * image_height), 
                        image_height - 1)
        landmark_point.append([landmark_x, landmark_y])
    return landmark_point

def pre_process_landmark(landmark_list):
    """Normalize landmarks to 0-1 range"""
    temp_landmark_list = copy.deepcopy(landmark_list)
    base_x, base_y = temp_landmark_list[0][0], temp_landmark_list[0][1]
    
    # Center on wrist (landmark 0)
    for i, point in enumerate(temp_landmark_list):
        temp_landmark_list[i][0] -= base_x
        temp_landmark_list[i][1] -= base_y
    
    # Flatten to 1D array (42 values = 21 landmarks × 2)
    temp_landmark_list = list(itertools.chain.from_iterable(
        temp_landmark_list))
    
    # Normalize by maximum absolute value
    max_value = max(list(map(abs, temp_landmark_list)))
    if max_value == 0:
        return temp_landmark_list
    
    normalized = [n / max_value for n in temp_landmark_list]
    return normalized
\end{lstlisting}

\paragraph{Preprocessing Strategy}
\begin{enumerate}
    \item \textbf{Wrist Centering}: Translate all points so wrist (index 0) is at origin
    \item \textbf{Flattening}: Convert 21 (x,y) pairs into 42-element vector
    \item \textbf{Normalization}: Divide by maximum value to scale to [-1, 1]
    \item \textbf{Invariance}: Makes gesture recognition scale and position invariant
\end{enumerate}

\subsection{Stage 4: Gesture Classification}

\paragraph{Keypoint Classifier Architecture}

\begin{lstlisting}[language=Python]
model = tf.keras.Sequential([
    tf.keras.layers.Input(shape=(42,)),              # 21 landmarks × 2
    tf.keras.layers.Dense(128, activation='relu'),   # Hidden layer 1
    tf.keras.layers.Dropout(0.2),                    # Regularization
    tf.keras.layers.Dense(64, activation='relu'),    # Hidden layer 2
    tf.keras.layers.Dropout(0.2),
    tf.keras.layers.Dense(num_classes, 
                         activation='softmax')        # Output layer
])

model.compile(
    optimizer=tf.keras.optimizers.Adam(learning_rate=0.001),
    loss='categorical_crossentropy',
    metrics=['accuracy']
)
\end{lstlisting}

\paragraph{Classification Process}

\begin{enumerate}
    \item Preprocessed landmarks (42 values) → Input layer
    \item Dense layer 1: 42 → 128 neurons with ReLU activation
    \item Dropout: Random 20\% deactivation to prevent overfitting
    \item Dense layer 2: 128 → 64 neurons with ReLU activation
    \item Dropout: 20\% to improve generalization
    \item Output layer: 64 → Number of gesture classes with Softmax
    \item Result: Probability distribution over gestures
\end{enumerate}

\paragraph{Decision Making}
\begin{lstlisting}[language=Python]
hand_sign_id = keypoint_classifier(pre_processed_landmark_list)
hand_sign_text = keypoint_classifier_labels[hand_sign_id]
confidence = max(classifier_output)

# Only accept if confidence > threshold
if confidence > 0.7:
    process_gesture(hand_sign_text)
\end{lstlisting}

\subsection{Stage 5: Point History Analysis (Complex Gestures)}

\paragraph{Purpose}

Recognize gestures involving motion (swipes, circles) using trajectory history:

\begin{lstlisting}[language=Python]
# Track finger tip (landmark 8 - index finger tip)
if hand_sign_id == 2:  # Point gesture
    point_history.append(landmark_list[8])
else:
    point_history.append([0, 0])

# Maintain history of last 16 frames
point_history = deque(maxlen=16)
\end{lstlisting}

\paragraph{Point History Preprocessing}
\begin{lstlisting}[language=Python]
def pre_process_point_history(image, point_history):
    """Normalize trajectory data"""
    image_width, image_height = image.shape[1], image.shape[0]
    temp_point_history = copy.deepcopy(point_history)
    
    base_x, base_y = temp_point_history[0][0], temp_point_history[0][1]
    
    for i, point in enumerate(temp_point_history):
        # Normalize by image dimensions
        temp_point_history[i][0] = (temp_point_history[i][0] - 
                                    base_x) / image_width
        temp_point_history[i][1] = (temp_point_history[i][1] - 
                                    base_y) / image_height
    
    # Flatten to 1D (32 values = 16 points × 2)
    temp_point_history = list(itertools.chain.from_iterable(
        temp_point_history))
    
    return temp_point_history
\end{lstlisting}

\paragraph{Output}
\begin{itemize}
    \item 32-element feature vector (16 frames × 2 coordinates)
    \item Trajectory-based gesture classification
    \item Used for detecting motion-based gestures
\end{itemize}

\subsection{Stage 6: Text Output Generation}

\paragraph{Label Lookup}
\begin{lstlisting}[language=Python]
with open('model/keypoint_classifier/keypoint_classifier_label.csv',
          encoding='utf-8-sig') as f:
    keypoint_classifier_labels = [row[0] for row in csv.reader(f)]

# Get text from prediction
gesture_id = keypoint_classifier(preprocessed_landmarks)
gesture_text = keypoint_classifier_labels[gesture_id]
\end{lstlisting}

\paragraph{Output Management}
\begin{itemize}
    \item Gesture text stored in variable
    \item Confidence score calculated
    \item Duplicate detection (skip if same as last)
    \item CSV logging for training data collection
\end{itemize}

\subsection{Stage 7: Speech Synthesis and Display}

\paragraph{Text-to-Speech Conversion}
\begin{lstlisting}[language=Python]
class SpeechEngine:
    def __init__(self):
        self._queue = queue.Queue()
        self._thread = threading.Thread(target=self._run, daemon=True)
        self._thread.start()

    def _run(self):
        engine = pyttsx3.init()
        while True:
            text = self._queue.get()
            if text is None:
                break
            engine.say(text)
            engine.runAndWait()

    def say(self, text: str):
        if text and pyttsx3 is not None:
            self._queue.put(text)

# Usage
speaker = SpeechEngine()
speaker.say("Hello")  # Non-blocking speech synthesis
\end{lstlisting}

\paragraph{Display Output}
\begin{lstlisting}[language=Python]
def draw_info_text(image, brect, handedness, hand_sign_text, 
                   finger_gesture_text):
    """Draw on image: hand type, gesture, finger gesture"""
    cv.rectangle(image, (brect[0], brect[1]), 
                 (brect[2], brect[1] - 22), (0, 0, 0), -1)
    
    info_text = handedness.classification[0].label[0:]
    if hand_sign_text != "":
        info_text = info_text + ':' + hand_sign_text
    
    cv.putText(image, info_text, (brect[0] + 5, brect[1] - 4),
               cv.FONT_HERSHEY_SIMPLEX, 0.6, (255, 255, 255), 1)
    
    return image
\end{lstlisting}

\section{Pipeline Flow Diagram}

\paragraph{Complete Data Flow}

\begin{verbatim}
┌──────────────────┐
│  Video Capture   │ (640×480, 30 FPS)
│   (Webcam)       │
└────────┬─────────┘
         │
         ▼
┌──────────────────┐
│  Preprocessing   │ (Flip, BGR→RGB, normalize)
│  (OpenCV)        │
└────────┬─────────┘
         │
         ▼
┌──────────────────┐
│ Hand Detection   │ (21 landmarks per hand)
│   (MediaPipe)    │
└────────┬─────────┘
         │
         ▼
┌──────────────────┐
│ Landmark         │ (Normalize, flatten to 42 values)
│ Preprocessing    │
└────────┬─────────┘
         │
    ┌────┴────┐
    │          │
    ▼          ▼
┌────────┐ ┌──────────────┐
│Keypoint│ │Point History │
│ Classify.│ │  Buffer      │
│(CNN)   │ │(16 frames)   │
└────┬───┘ └──────┬───────┘
     │             │
     └─────┬───────┘
           │
           ▼
    ┌────────────────┐
    │ Label Lookup   │ (Get gesture name)
    └────────┬───────┘
             │
             ▼
    ┌─────────────────┐
    │ Text Output     │ (Display + CSV log)
    └────────┬────────┘
             │
             ▼
    ┌─────────────────┐
    │Text-to-Speech   │ (pyttsx3, threading)
    │Engine           │
    └────────┬────────┘
             │
             ▼
    ┌─────────────────┐
    │ Display Results │ (UI, landmarks, speech)
    │ (OpenCV/Streamlit)│
    └─────────────────┘
\end{verbatim}

\section{Key Pipeline Features}

\subsection{Dual-Hand Support}

Changed configuration from single to dual-hand recognition:
\begin{lstlisting}[language=Python]
# Before: num_hands=1
# After:  num_hands=2

options = vision.HandLandmarkerOptions(
    num_hands=2,  # Support both hands simultaneously
    ...
)
\end{lstlisting}

Each hand is processed independently through the complete pipeline.

\subsection{Real-Time Processing}

Pipeline optimizations for real-time operation:
\begin{itemize}
    \item FPS calculation and monitoring
    \item Video mode for temporal consistency
    \item Threading for non-blocking operations
    \item Efficient numpy operations for mathematics
\end{itemize}

\subsection{Data Augmentation}

During dataset collection, automatically achieved through:
\begin{itemize}
    \item Varying hand distances from camera
    \item Different lighting conditions
    \item Multiple hand positions in frame
    \item Both hands and hand angles
\end{itemize}

This natural augmentation ensures robustness without explicit transformation functions.
